% Fixture 5: notes via the pdfcomment package, in several annotation shapes.
% Frame 5 carries a FreeText annotation that IS visible content and must not be treated as a note.
\documentclass[aspectratio=169]{beamer}
\usetheme{default}
\usepackage{pdfcomment}
\title{pdfcomment notes}
\subtitle{Fixture 5}
\author{Lectern fixtures}
\date{}
\begin{document}
\begin{frame}
  \titlepage
  \pdfcomment{Title page note via pdfcomment with the default icon.}
\end{frame}
\section{Main}
\begin{frame}{Note icon with subject}
  Slide content.
  \pdfcomment[icon=Note,color=yellow,subject={Speaker notes}]{A note with the Note icon, a colour and a subject.}
\end{frame}
\begin{frame}{Margin note with overlays}
  \pdfmargincomment{A margin note (pdfmargincomment).}
  Build one.
  \pause
  Build two.
\end{frame}
\begin{frame}{Open comment with author}
  Slide text.
  \pdfcomment[open=true,author={Ben},icon=Insert]{A note whose popup is open (/Open true) and that carries an author (/T).}
\end{frame}
\begin{frame}{Free text (visible annotation)}
  Slide text.
  \pdffreetextcomment[width=5cm,height=1cm,voffset=-2cm,color=red,fontcolor=black]{A FreeText annotation: visible content, not a note.}
\end{frame}
\begin{frame}{No annotations}
  Nothing attached.
\end{frame}
\end{document}
